\section{Axis 2: Topological Formalisms}
\label{sec:topology}

Topology enters granular computing through two research communities that
develop largely independently.  We call them the \emph{abstract topology}
branch and the \emph{computational topology (TDA)} branch.  Mapping this
split is one of the original contributions of this survey.

We state the separation as a bounded, checkable claim rather than as a
universal one.  Within the corpus assembled by the search protocol of
Appendix~\ref{app:search}, no paper cites work from both branches, and we
found none in either branch's reference list that reaches into the other.
This is a statement about a 52-paper corpus retrieved from OpenAlex and
Google Scholar in September 2025, not a proof that no such paper exists
anywhere: the two branches publish in disjoint venue sets, and a bridging
paper in a venue outside our search frame would not have been retrieved.
What the corpus does establish is that the bridge is rare enough to be
absent from the literature a researcher entering this area would find.

%-----------------------------------------------------------------------
\subsection{Branch A: Abstract Topology}
\label{sec:abstract_topo}
%-----------------------------------------------------------------------

Abstract topology treats the neighbourhood operator $\NB{B}$ as a
topological structure and asks: which axioms does it satisfy?  What is the
algebraic relationship between the lower/upper approximation operators and
classical topological operators (interior, closure)?  When do two different
neighbourhood systems produce the same approximation operators?

\paragraph{Foundational work.}
Zhu~\cite{Zhu2007covering} established covering rough sets from the topological
viewpoint: a \emph{covering} of $\U$ generates a neighbourhood system; the
paper characterises when two coverings share the same lower (or upper)
approximation operator, and constructs complete axiomatic systems for both.
With 637 citations, this is the most-cited Axis~2 paper in this survey.
D'eer et al.~\cite{DeerCornelisYao2016} later showed that several
covering-based rough set models in the literature are semantically unsound---they
fail to reduce to Pawlak rough sets when the covering is a partition---and
proposed corrected definitions that preserve the algebraic link to the
original framework.

Yao~\cite{Yao2023Alexandrov} identifies the Alexandrov topology as the
unifying structure: a preorder $\preceq$ on $\U$ induces an Alexandrov
topology whose specialisation preorder recovers $\preceq$.  Under this lens,
the four main rough set models (binary-relation, neighbourhood-operator,
covering, subset-system) are all instances of the same topological framework,
distinguished only by which axioms (reflexivity, transitivity, symmetry) the
preorder satisfies.

\paragraph{Separation axioms.}
A parallel line studies which topological separation axioms (T$_0$ through
T$_4$) are satisfied by the approximation space induced by an equivalence or
neighbourhood relation, and whether stronger axioms correspond to better
approximation quality (smaller boundary region).  The common thread across
four recent papers---Al-shami et al.~\cite{AlShami2025delta} ($\delta$-open
sets), Al-shami~\cite{AlShami2022topology} (somewhat-open sets via
$N_\rho$-neighbourhoods), Al-shami and Alshammari~\cite{AlShami2022supra}
(supra-topology), and El~Safty~\cite{ElSafty2021}
(simply$^*\alpha$-open sets)---is the progressive weakening of the open-set
axiom: each relaxation yields coarser granules and therefore a smaller
boundary region, at the cost of weaker algebraic properties.  None of
these papers benchmarks the resulting approximation against NRS/GBNRS at
the scale of Section~\ref{sec:methodology}'s recommended protocol, so
the practical magnitude of the boundary-region improvement remains
unquantified.

\paragraph{Fuzzy topology.}
Lai and Zhang's~\cite{LaiZhang2006} equivalence between fuzzy preorders and
fuzzy Alexandrov topologies provides a bridge to fuzzy rough sets: an $\alpha$-cut
of a fuzzy similarity relation $R$ yields a crisp partition at threshold
$\alpha$, and the resulting nested family of partitions defines a fuzzy
topology on $\U$.  The authors' prior work~\cite{DaiTT2024IFT} exploits
this structure by building an \emph{intuitionistic} fuzzy topology (IFT)
from a pre-order relation, then using the IFT subbasis to measure attribute
significance; the paper reports reducts 50\% smaller on average than
dependency-based NRS, with 10\% higher accuracy.
Applying this survey's own diagnostic: the IFT subbasis significance measure
is a monotone function of the pre-order's $\alpha$-cut chain---placing it in
the Pattern~A risk zone (reduction to NRS dependency at the cut-threshold).
The 50\%/10\% figures are also reported without significance testing, raising
the Pattern~C concern; we flag these as claims requiring independent
verification rather than taking them at face value.

\paragraph{Assessment of the abstract branch.}
The abstract topology branch is mathematically mature.  Its primary
contribution has been taxonomic: clarifying which set-theoretic properties
of an approximation operator correspond to which topological axioms.  The
computational payoff---smaller reducts, better accuracy---is reported in
isolated papers but has not been benchmarked systematically against modern
NRS/GBNRS at the scale and with the rigour advocated in
Section~\ref{sec:methodology}.

%-----------------------------------------------------------------------
\subsection{Branch B: Computational Topology (TDA)}
\label{sec:tda}
%-----------------------------------------------------------------------

The TDA branch uses \emph{persistent homology} (Section~\ref{sec:ph}) as
a computational tool and asks: can the topological summary of a dataset
generate neighbourhoods, identify important attributes, or define a rough
approximation space?

\paragraph{Background.}
Carlsson~\cite{Carlsson2009topology} established persistent homology as a
tool for data analysis, showing that the birth/death intervals of topological
features across scale form a stable, discriminative summary;
Ghrist~\cite{Ghrist2008barcodes} introduced the \emph{barcode} representation
that encodes these intervals compactly.
Chazal and Michel~\cite{ChazalMichel2021TDA} provide a practical introduction
covering filtrations, persistence diagrams, and vectorisation for practitioners;
Otter et al.~\cite{Otter2017roadmap} survey the computational algorithms
for building different filtration types.
Hensel et al.~\cite{HenselMoorRieck2021} survey topological machine learning
methods beyond point-cloud homology, including graph persistence and topological
regularisers in deep learning.  None of these methods interact with rough set
approximation operators.
Pun et al.~\cite{Pun2022PH} survey persistent-homology-based machine learning
across classification, regression, and feature-extraction tasks, and provide
the only systematic comparison of simplicial complex types (alpha vs.\
Vietoris--Rips) and vectorisation strategies.  None of the 50+ methods they
survey are applied to rough-set reducts; this is one of the two cross-branch
gaps this survey maps.

\paragraph{Classifier via simplicial granules.}
Kindelan et al.~\cite{Kindelan2021TDABC} build a filtered simplicial complex
on the dataset, use persistent homology to select a sub-complex that captures
the data's topology, and define each unlabelled point's neighbourhood via its
\emph{star} and \emph{link} operators in the sub-complex.  Labels propagate
from labelled to unlabelled points through these simplicial neighbourhoods.
On eight datasets including imbalanced and overlapping classes, the method
outperforms $k$-NN and weighted $k$-NN.

This is, in the language of granular computing, a method for defining an
\emph{adaptive granule} $\NB{}(x) = \mathrm{star}(x) \cup \mathrm{link}(x)$
from the dataset's topology.  Kindelan et al.\ do not use the language of
rough sets and do not compute reducts or boundary regions; their goal is
classification, not feature selection.

\paragraph{Attribute importance from persistence.}
Our own work (Sections~\ref{sec:patB}--\ref{sec:patC}) tests whether
persistence diagrams can score attribute importance.  Writing
$\mathrm{Dgm}_k(\U_B)$ for the dimension-$k$ persistence diagram of the
Vietoris--Rips filtration built on the projection of $\U$ onto attribute
set $B$, we define the \emph{persistence attribute importance} of $a_j$ as
the total bottleneck displacement caused by deleting that coordinate,
summed over the homological dimensions actually computed:
\begin{equation}\label{eq:pai_def}
  \mathrm{PAI}(j) \;=\; \sum_{k=0}^{d_{\max}}
    d_B\!\left(\mathrm{Dgm}_k(\U_{\C}),\,
               \mathrm{Dgm}_k(\U_{\C\setminus\{j\}})\right),
\end{equation}
with $d_{\max}=1$ throughout (so $H_0$ and $H_1$ contribute).  The sum,
rather than a single diagram distance, is what our implementation computes,
and it is what the following bound is stated for.

\begin{theorem}[PAI upper bound]
\label{thm:pai_bound}
Let $\U$ be a finite point cloud, let all attributes be real-valued, and let
the Vietoris--Rips filtrations be built from the Euclidean metric on the
retained coordinates.  Then
\begin{equation}\label{eq:pai_bound}
  \mathrm{PAI}(j) \leq (d_{\max}+1)\,\mathrm{range}(a_j),
\end{equation}
where $\mathrm{range}(a_j) = \max_{x\in\U} a_j(x) - \min_{x\in\U} a_j(x)$.
Moreover each individual summand in~\eqref{eq:pai_def} obeys the sharper
bound $d_B(\mathrm{Dgm}_k(\U_{\C}),\mathrm{Dgm}_k(\U_{\C\setminus\{j\}}))
\leq \mathrm{range}(a_j)$.
\end{theorem}

\begin{proof}
Let $d_{\C}$ and $d_{\C\setminus\{j\}}$ denote the pairwise Euclidean distance
functions on $\U$ induced by attribute sets $\C$ and $\C\setminus\{j\}$.  For
any $x,y\in\U$, writing $S=\sum_{a_i\in\C\setminus\{j\}}(a_i(x)-a_i(y))^2$ and
$t=(a_j(x)-a_j(y))^2$, subadditivity of the square root gives
$\sqrt{S+t}-\sqrt{S}\leq\sqrt{t}$, so
\[
  0 \;\leq\; d_{\C}(x,y) - d_{\C\setminus\{j\}}(x,y)
    \;\leq\; |a_j(x)-a_j(y)| \;\leq\; \mathrm{range}(a_j),
\]
hence $\|d_{\C}-d_{\C\setminus\{j\}}\|_\infty \leq \mathrm{range}(a_j)$.
Two Vietoris--Rips filtrations whose defining distance functions differ by at
most $\varepsilon$ in the supremum norm are $\varepsilon$-interleaved, so by
the algebraic stability theorem for persistence
modules~\cite{CohenSteiner2007stability,Chazal2009Stability} each dimension
satisfies
$d_B(\mathrm{Dgm}_k(\U_{\C}),\mathrm{Dgm}_k(\U_{\C\setminus\{j\}}))
\leq \|d_{\C}-d_{\C\setminus\{j\}}\|_\infty \leq \mathrm{range}(a_j)$.
This is the sharper claim.  Summing it over the $d_{\max}+1$ dimensions
in~\eqref{eq:pai_def} yields~\eqref{eq:pai_bound}.
\end{proof}

The multiplier $(d_{\max}+1)$ is therefore not a property of the Rips
skeleton but simply a count of the terms being summed; a single-dimension
importance score would carry no multiplier.  The released column
\texttt{pai\_bound} in \texttt{raw\_importance.csv} records the right-hand
side of~\eqref{eq:pai_bound}, i.e.\ $2\,\mathrm{range}(a_j)$ at
$d_{\max}=1$.

\begin{remark}[Domain of validity]
\label{rem:pai_domain}
Equation~\eqref{eq:pai_def} presupposes $|\C|\geq 2$: at $d=1$ the reduced
cloud $\U_{\C\setminus\{j\}}$ is empty and PAI is undefined.  The scores are
also invariant to translation of $a_j$ but not to rescaling, which is the
formal source of the measurement-scale special case of Pattern~B
(Section~\ref{sec:patB}).
\end{remark}

The bound is valid but loose: across all ten datasets and five seeds there
are zero violations, with mean ratio
$\mathrm{PAI}(j)/[(d_{\max}+1)\mathrm{range}(a_j)] = 0.038$ (max $0.194$);
against the sharper per-dimension bound the mean ratio is $0.075$
(max $0.387$).  Because it is loose by more than an order of magnitude, the
bound does not by itself imply that PAI tracks range---it is consistent with
that hypothesis and explains why the hypothesis is worth testing.  What the
bound does establish is structural: PAI's admissible magnitude is
proportional to the feature's range, so a feature can score highly for
reasons unrelated to class structure.  Whether range is in fact the dominant
driver is an empirical question, tested directly in Section~\ref{sec:patB}.

\paragraph{Alternative vectorisation strategies.}
The TDA literature offers several vectorisation approaches beyond the
bottleneck distance used in this survey's experiments.  Persistence
landscapes~\cite{Bubenik2015landscapes} map each diagram to a sequence of
piecewise-linear functions amenable to Banach-space statistics; persistence
images~\cite{Adams2017persistence} discretise the birth-death plane into a
stable, fixed-dimensional vector; and the Mapper
algorithm~\cite{Singh2007mapper} produces a simplicial summary of the data's
shape without computing homology.  Pun et al.~\cite{Pun2022PH} survey these
and 50+ further methods.  We chose bottleneck distance because it directly
pairs with the stability theorem underlying Theorem~\ref{thm:pai_bound};
whether alternative vectorisations yield importance scores that correlate
more strongly with class structure is an open question.

\paragraph{Broader TDA survey.}
Su et al.~\cite{Su2025TDA} survey TDA beyond persistent homology, covering
topological Laplacians, Dirac operators, sheaf theory, and Hodge
decomposition.  None of these richer tools have been applied to granular
computing's core pipeline; this remains an open frontier, though
the risk of Pattern~A (reduction to a known granular quantity) remains.

%-----------------------------------------------------------------------
\subsection{The Gap Between the Two Branches}
\label{sec:topo_gap}
%-----------------------------------------------------------------------

The abstract topology branch (Branch~A) and the computational topology
branch (Branch~B) have evolved in parallel without cross-pollination.
Figure~\ref{fig:topo_gap} visualises this split: within our corpus, no paper
jointly defines a rough approximation space using persistent-homological
granules and characterises the separation axioms of that space.

\begin{figure}[pos=!htbp]
\centering
\begin{tikzpicture}[
  branchbox/.style = {draw, rounded corners=5pt, fill=gray!8,
                      minimum width=3.6cm, minimum height=4.8cm,
                      align=center},
  paper/.style     = {draw, rounded corners=3pt,
                      minimum width=3.2cm, minimum height=0.55cm,
                      align=center, font=\scriptsize, fill=white},
  paperA/.style    = {paper, fill=blue!10,   draw=blue!50!black},
  paperB/.style    = {paper, fill=green!10,  draw=green!55!black},
  titlestyle/.style= {font=\small\bfseries},
  arr/.style       = {-{Stealth[length=4pt]}, thick},
  gaparr/.style    = {<->, >=Stealth, thick, dashed, red!60!black},
]

%% ── Branch A column ──────────────────────────────────────────────────────────
\node[branchbox] (boxA) at (0,0) {};
\node[titlestyle, above=3pt] at (boxA.north) {Branch A};
\node[font=\scriptsize\itshape, below=1pt] at (boxA.north) {Abstract Topology};

\node[paperA] (zA)  at (0, 1.5) {Zhu 2007 (covering RS)};
\node[paperA] (yA)  at (0, 0.8) {Yao 2023 (Alexandrov)};
\node[paperA] (aA)  at (0, 0.1) {Al-shami 2022/2025};
\node[paperA] (dA)  at (0,-0.6) {D'eer et al.\ 2016};
\node[paperA] (lA)  at (0,-1.3) {Lai \& Zhang 2006};

\node[font=\scriptsize, align=center] at (0,-2.05)
  {\textit{Separation axioms,}\\[-2pt]\textit{covering properties,}\\[-2pt]\textit{axiomatic systems}};

%% ── Branch B column  (x = 7.0, gap = 3.4 cm) ────────────────────────────────
\node[branchbox] (boxB) at (7.0,0) {};
\node[titlestyle, above=3pt] at (boxB.north) {Branch B};
\node[font=\scriptsize\itshape, below=1pt] at (boxB.north) {Computational TDA};

\node[paperB] (kB)  at (7.0, 1.5) {Kindelan 2021 (TDABC)};
\node[paperB] (sB)  at (7.0, 0.8) {Su 2025 (TDA survey)};
\node[paperB] (dB)  at (7.0, 0.1) {Dai 2024 (IFT)};
\node[paperB] (pB)  at (7.0,-0.6) {Pun 2022 (PH-ML)};
\node[paperB] (hB)  at (7.0,-1.3) {Hensel et al.\ 2021};

\node[font=\scriptsize, align=center] at (7.0,-2.05)
  {\textit{Persistence diagrams,}\\[-2pt]\textit{filtrations, vectorisation,}\\[-2pt]\textit{classification}};

%% ── Zero-link zone in between  (centre of gap = 3.5) ────────────────────────
\draw[gaparr] (boxA.east) -- node[above, font=\scriptsize\bfseries, red!60!black]
  {0 joint papers in corpus} (boxB.west);

\node[font=\scriptsize, align=center, text=red!60!black]
  at (3.5,-0.3) {\textit{Missing:}\\[-2pt]\textit{rough approx.\ space}\\[-2pt]\textit{from PH granules}\\[-2pt]\textit{+ separation axioms}};

%% ── Background shading ───────────────────────────────────────────────────────
\begin{scope}[on background layer]
  \node[fill=blue!4, rounded corners=6pt,
        fit=(boxA), inner sep=2pt] {};
  \node[fill=green!4, rounded corners=6pt,
        fit=(boxB), inner sep=2pt] {};
\end{scope}

\end{tikzpicture}
\caption{The two topology branches in this survey are not connected by any
paper in our corpus (Appendix~\ref{app:search}; retrieved September 2025).
Branch~A (abstract topology) establishes which algebraic
axioms GrC approximation spaces satisfy; Branch~B (computational TDA) uses
persistent homology for data analysis.  The gap is the missing theorem that
connects persistence-based granules to separation axioms --- an open problem
identified in Section~\ref{sec:open}.}
\label{fig:topo_gap}
\end{figure}

This gap is not merely bibliographic.  Branch~A shows that covering axioms
constrain which concepts can be exactly approximated; Branch~B provides a
data-driven mechanism for building coverings from geometric-topological
summaries.  A paper connecting the two---showing, for example, that
persistent-homological granules satisfy a T$_1$ or T$_2$ separation axiom
under certain dataset conditions---would be a new result.

However, Pattern~A warns that such a connection must be checked carefully:
the star operator in a Vietoris--Rips complex at scale $\varepsilon$ is,
for finite point clouds, equivalent to the fixed-radius ball neighbourhood
$\NB{}^\varepsilon$ when restricted to the input points.  A naive connection
reduces to NRS.  The contribution must come from using the \emph{persistent
filtration}---the birth/death structure across scales---rather than any
single scale.

\paragraph{Practical guidance.}
Any researcher wishing to bridge Branch~A and Branch~B should:
(i) define the persistence-based granule explicitly and check
    against Definitions~\ref{def:nb}--\ref{def:dep};
(ii) apply the Pattern~A diagnostic (Section~\ref{sec:patA}): does the
     granule at a fixed filtration scale reduce to $\delta$-ball NRS?  If
     yes, the contribution must derive from the multi-scale structure;
(iii) verify experimentally that the resulting reduct is smaller or its
     boundary region is smaller than GBNRS on the same benchmark (a high bar
     given GBNRS's 259 citations and well-tuned defaults).
